\subsection*{一步一步建立直观}

\paragraph{1. 把复合函数拆成可求导的环节}

记 $z=w_1x_1+w_2x_2+b$，输出为 $a=\max(0,z)$，损失为 $L=(a-y)^2$。对 $w_1$ 求导时，依次经过 $L$、$a$、$z$ 三个环节。右图同时给出中间量，方便逐项核对链式法则中的每个因子，而不是只看最后一个梯度数字。每一环的输入位置都来自同一次前向计算。



\paragraph{2. 用符号解释参数应该向哪边走}

在激活区域，权重梯度为 $2(a-y)x_1$。如果输出过大且输入为正，梯度为正，梯度下降就减小这个权重；若输入为负，改变方向随之反转。这种符号推理是检查求导的有效方法。偏置对应的局部输入恒为一，所以少了输入因子。



\paragraph{3. 不可导点需要单独处理}

把单元切到关闭侧，再靠近边界。严格负侧梯度为零，严格正侧使用前面的公式；边界两侧一般不一致。右图在零点明确提示折点，不给出通常导数；某些训练程序会约定传回零，但这不是证明 ReLU 在零点可导。形式化命题与手算推导的适用条件应按各自陈述逐一核对。



\subsection*{定义、命题与证明}

\paragraph{定义：单个 ReLU 单元与平方损失}

记 $z=w_1x_1+w_2x_2+b$、$a=\max(0,z)$、$L=(a-y)^2$。输入与目标固定时，损失是参数的复合函数。

\paragraph{命题：严格激活区域中的权重梯度}

在 $z>0$ 的邻域中，$\partial L/\partial w_1=2(a-y)x_1$。在 $z<0$ 的邻域中，该梯度为零。

\begin{proof}

当 $z>0$ 时，ReLU 在邻域内就是恒等函数，链式法则给出 $2(a-y)\cdot1\cdot x_1$。当 $z<0$ 时，它在邻域内恒为零，局部损失因而不随该参数变化。

结论刻意避开 $z=0$。ReLU 在零处没有经典导数，复合损失是否恰好可导还要另看具体参数与目标；程序采用的反向传播约定不能替代这项数学判断。

\end{proof}
